% Gesamttest «Trigonometrische Funktionen» — Quelle des PDFs (python3 scripts/build-lp-pdf.py trigonometrische).
% Musterlösung und Raster: bewertungspaket.tex. Alle Werte mit python3 nachgerechnet (05.10.2026).
% Kein \lt/\gt — pdfLaTeX kennt sie nicht, < und > tun es.
\documentclass[11pt]{article}
\usepackage{../lp-druck}
\renewcommand{\lpname}{Trigonometrische Funktionen}
\renewcommand{\lpteilgebiet}{3.5}
\renewcommand{\lpbereich}{Schwerpunktfach}
\renewcommand{\lpfuss}{Gesamttest · 25 Punkte}
\begin{document}
\lpkopf{Gesamttest}{%
  \vspace{4pt}\noindent\begin{tabularx}{\linewidth}{@{}X X@{}}
  Code (kein Name): \hrulefill & Punkte: \hrulefill\ / 25
  \end{tabularx}}

\begin{lpinfo}{Anleitung}
\textbf{Zeit:} rund 30 Minuten. \textbf{Hilfsmittel:} Die Kompetenz des Teilgebiets 3.5 trägt im Lehrplan den Vermerk «mit und ohne Hilfsmittel». Darum: \textbf{Teil A ohne Taschenrechner}, erst für \textbf{Teil B} den Taschenrechner nehmen (im Bogenmass, RAD). Lineal und Bleistift sind erlaubt.
Schreib zu jeder Aufgabe den \textbf{Rechenweg} auf — bewertet wird der Weg.
Danach: Lösung scannen oder fotografieren (ohne Namen und Standort) und mit dem \emph{Bewertungspaket} bewerten lassen.
\end{lpinfo}

\lpteil{Teil A · ohne Taschenrechner}{Kapitel 1 bis 4 · 15 Punkte}

\begin{lpaufgabe}{G1}{Sinus und Cosinus skizzieren}{3 P}
Skizziere $y = \sin x$ und $y = \cos x$ für $-\pi \le x \le 2\pi$ in das Koordinatensystem. Gib alle Nullstellen von $y = \cos x$ in diesem Bereich an.
\par\smallskip\centering
\begin{tikzpicture}
\begin{axis}[width=15.5cm,height=6.2cm,axis lines=middle,xmin=-3.5,xmax=6.9,ymin=-1.5,ymax=1.5,
  xtick={-3.14159,-1.5708,1.5708,3.14159,4.71239,6.28319},xticklabels={$-\pi$,$-\frac{\pi}{2}$,$\frac{\pi}{2}$,$\pi$,$\frac{3\pi}{2}$,$2\pi$},
  ytick={-1,1},minor ytick={-0.5,0.5},minor xtick={-2.35619,-0.785398,0.785398,2.35619,3.92699,5.49779},grid=both,grid style={lplinie},tick label style={font=\scriptsize},
  xlabel=$x$,ylabel=$y$,
  yticklabel style={fill=white,inner sep=1pt,xshift=-2pt},xticklabel style={fill=white,inner sep=1pt,yshift=-2pt}]
\end{axis}
\end{tikzpicture}
\par\raggedright
\lpplatz{2}
\end{lpaufgabe}

\begin{lpaufgabe}{G2}{Werte ohne Rechner}{3 P}
\begin{enumerate}[label=(\alph*),leftmargin=*,itemsep=1pt]
\item Rechne $\tfrac{5\pi}{3}$ in Grad um.
\item Gib an: $\sin \tfrac{11\pi}{6}$ und $\cos \tfrac{4\pi}{3}$.
\item Gib an: $\tan \tfrac{3\pi}{2}$ und $\tan\left(-\tfrac{3\pi}{4}\right)$.
\end{enumerate}
\lpplatz{5}
\end{lpaufgabe}

\begin{lpaufgabe}{G3}{Symmetrie und Periode}{3 P}
Es gilt $\sin 1.1 \approx 0.891$ und $\cos 1.1 \approx 0.454$. Gib ohne Taschenrechner an und begründe jeweils kurz:
\begin{enumerate}[label=(\alph*),leftmargin=*,itemsep=1pt]
\item $\sin(-1.1)$
\item $\cos(-1.1)$
\item $\sin(\pi - 1.1)$
\end{enumerate}
\lpplatz{5}
\end{lpaufgabe}

\begin{lpaufgabe}{G4}{Die Tangensfunktion}{3 P}
\begin{enumerate}[label=(\alph*),leftmargin=*,itemsep=1pt]
\item Gib die Definitionsmenge und die Periodenlänge von $y = \tan x$ an.
\item Gib alle Polstellen und alle Nullstellen von $y = \tan x$ im Intervall $[0;\, 2\pi]$ an.
\end{enumerate}
\lpplatz{5}
\end{lpaufgabe}

\begin{lpaufgabe}{G5}{Kenngrössen}{3 P}
Gegeben ist $y = 2\sin(3x) - 1$. Gib Amplitude, Periodenlänge und Wertemenge an.
\lpplatz{4}
\end{lpaufgabe}

\lpteil{Teil B · mit Taschenrechner}{Kapitel 4 und 5 · 10 Punkte}

\begin{lpaufgabe}{G6}{Gleichung aus dem Graphen}{3 P}
Die Kurve hat die Form $y = a \sin(bx) + v$ mit $a > 0$. Bestimme $a$, $b$ und $v$ und begründe mit Werten aus dem Graphen.
\par\smallskip\centering
\begin{tikzpicture}
\begin{axis}[width=15.5cm,height=6.6cm,axis lines=middle,xmin=-0.6,xmax=13.2,ymin=-2.4,ymax=4.4,
  xtick={3.14159,6.28319,9.42478,12.56637},xticklabels={$\pi$,$2\pi$,$3\pi$,$4\pi$},
  minor xtick={1.5708,4.71239,7.85398,10.99557},
  ytick={-2,-1,1,2,3,4},grid=both,grid style={lplinie},tick label style={font=\scriptsize},
  xlabel=$x$,ylabel=$y$,
  yticklabel style={fill=white,inner sep=1pt,xshift=-2pt},xticklabel style={fill=white,inner sep=1pt,yshift=-2pt}]
\addplot[lpblau,very thick,domain=-0.6:13.2,samples=200]{2.5*sin(deg(x/2))+1};
\end{axis}
\end{tikzpicture}
\par\raggedright
\lpplatz{2}
\end{lpaufgabe}

\begin{lpaufgabe}{G7}{Alle Lösungen}{4 P}
Löse im Intervall $[0;\, 2\pi]$, auf drei Dezimalen:
\begin{enumerate}[label=(\alph*),leftmargin=*,itemsep=1pt]
\item $\sin x = 0.35$
\item $\cos x = 0.35$
\end{enumerate}
\lpplatz{4}
\end{lpaufgabe}

\begin{lpaufgabe}{G8}{Ebbe und Flut}{3 P}
Der Wasserstand in einem Hafen ist $h(t) = 2\sin\left(\tfrac{\pi}{6}\,t\right) + 3$ ($h$ in m, $t$ in Stunden ab Mitternacht).
\begin{enumerate}[label=(\alph*),leftmargin=*,itemsep=1pt]
\item Wie hoch steht das Wasser höchstens, wie hoch mindestens?
\item Nach wie vielen Stunden wiederholt sich der Verlauf?
\item Zu welchen Zeiten zwischen $0$ und $12$ Uhr steht das Wasser genau $4$ m hoch?
\end{enumerate}
\lpplatz{4}
\end{lpaufgabe}
\end{document}
